\section*{Chapitre \ref{ch:g11:plate-motions} --- Les plaques et leurs mouvements}

\begin{solution}{exo:g11:plate-motions:1}
Une \omterm{def:g11:plate-motions:rigid-plate}{plaque rigide} est un morceau de \omterm{def:g11:lithosphere-asthenosphere:lithosphere}{lithosphère} qui se déplace d'un seul
bloc~: les distances entre ses points ne changent pas, et sa déformation est
concentrée dans des zones étroites le long de ses limites. Les séismes le
montrent~: presque tous se produisent le long de ceintures étroites, les
limites de plaques, et très peu à l'intérieur des plaques.
\end{solution}

\begin{solution}{exo:g11:plate-motions:2}
\qty{72}{mm} par an correspondent à \qty{72}{km} par million d'années. En
30 ans, la station se déplace de $72 \times 30 = \qty{2160}{mm}$, environ
\qty{2.2}{m}~; en un million d'années, de \qty{72}{km}.
\end{solution}

\begin{solution}{exo:g11:plate-motions:3}
$v_E = 30/10 = \qty{3}{cm}$ par an et $v_N = 40/10 = \qty{4}{cm}$ par an. La
vitesse vaut $\sqrt{3^2 + 4^2} = \qty{5}{cm}$ par an. L'angle $\beta$ avec
le nord vérifie $\tan\beta = 3/4$, donc $\beta \approx \ang{37}$~: la
station se déplace vers le nord-est, à \ang{37} à l'est du nord.
\end{solution}

\begin{solution}{exo:g11:plate-motions:4}
Une \omterm{def:g11:plate-motions:hotspot-track}{trace de point chaud} est l'alignement de volcans et de monts sous-marins
qu'une plaque emporte loin d'un point chaud en se déplaçant au-dessus de
lui, de plus en plus anciens à mesure qu'on s'éloigne du volcan actif. Le
volcan le plus jeune, Kilauea, se trouve à l'extrémité sud-est de la chaîne,
sur l'île d'Hawaï~; les âges augmentent vers l'ouest-nord-ouest. La plaque
emporte ses volcans loin du point chaud dans la direction de son
mouvement~: la plaque Pacifique se déplace vers l'ouest-nord-ouest.
\end{solution}

\begin{solution}{exo:g11:plate-motions:5}
$519/5.1 \approx \qty{102}{km}$ par million d'années, soit \qty{102}{mm} par
an, environ \qty{10}{cm} par an.
\end{solution}

\begin{solution}{exo:g11:plate-motions:6}
$v_E = 63/31 \approx \qty{20.3}{mm}$ par an et $v_N = 48/31 \approx
\qty{15.5}{mm}$ par an. La vitesse vaut $\sqrt{20.3^2 + 15.5^2} \approx
\qty{25.5}{mm}$ par an~; $\tan\alpha = 15.5/20.3 \approx 0.76$ donne $\alpha
\approx \ang{37}$~: la station se déplace à \ang{37} au nord de l'est, vers
le nord-est. Le Mauna Kea se déplace $72/25.5 \approx 2.8$ fois plus vite.
\end{solution}

\begin{solution}{exo:g11:plate-motions:7}
(1) En millimètres par an,
\[
  \vec v_{\mathrm{NA/EU}} = (-10.5 - 7.7,\ 20.6 - 16.3) = (-18.2, +4.3).
\] Sa norme vaut $\sqrt{18.2^2 + 4.3^2} \approx
\qty{18.7}{mm}$ par an, et $\tan\alpha = 4.3/18.2 \approx 0.24$ donne
\ang{13} au nord de l'ouest. (2) Les deux plaques s'écartent de part et
d'autre de la dorsale d'environ \qty{1.9}{cm} par an~: c'est la \omterm{def:g11:drift-to-spreading:spreading-rate}{vitesse
d'expansion totale} de la dorsale en Islande. (3) Elle concorde avec les
\qty{2}{cm} par an cités pour l'Atlantique Nord.
\end{solution}

\begin{solution}{exo:g11:plate-motions:8}
(1) $v = \pi \times 0.651 \times 6371/180 \approx \qty{72.4}{km}$ par million
d'années, environ \qty{7.2}{cm} par an. (2) $72.4 \times \sin\ang{30} \approx
\qty{36.2}{km}$ par million d'années, environ \qty{3.6}{cm} par an. (3) Au
pôle de rotation, s'il se trouve sur la plaque. La vitesse change d'un point
à l'autre, de zéro au pôle à sa valeur maximale à \ang{90} de celui-ci~: une
plaque entière n'a pas de vitesse unique, seulement une vitesse en chaque
point (ou une moyenne sur sa surface).
\end{solution}

\begin{solution}{exo:g11:plate-motions:9}
(1) En millimètres par an,
\[
  \vec v_{\mathrm{IN/EU}} = (40.6 - 24.0,\ 38.6 + 3.3) = (+16.6, +41.9).
\] La vitesse vaut $\sqrt{16.6^2 + 41.9^2} \approx
\qty{45}{mm}$ par an~; l'angle $\beta$ avec le nord vérifie $\tan\beta =
16.6/41.9 \approx 0.40$, donc $\beta \approx \ang{22}$~: l'Inde se rapproche
de l'Eurasie, à \ang{22} à l'est du nord. (2) \qty{45}{km} par million
d'années, donc \qty{450}{km} en 10 millions d'années.
\end{solution}

\begin{solution}{exo:g11:plate-motions:10}
$(2432 - 1058)/(27.7 - 10.3) = 1374/17.4 \approx \qty{79}{km}$ par million
d'années, environ \qty{7.9}{cm} par an, un peu moins que \qty{9.1}{cm} par
an. Les âges sont incertains d'environ un demi-million d'années chacun~; la
différence d'âge, $17.4 \pm 1.0$ millions d'années, donne une vitesse
comprise entre environ 75 et \qty{84}{km} par million d'années. Entre ces
deux volcans, la plaque semble s'être déplacée un peu plus lentement qu'en
moyenne.
\end{solution}

\begin{solution}{exo:g11:plate-motions:11}
(1) $\vec v_{A/C} = \vec v_{A/B} + \vec v_{B/C}$~: \qty{3}{cm} par an vers
l'est plus \qty{4}{cm} par an vers le nord. Les deux flèches forment un
angle droit, et le troisième côté du triangle a pour longueur
$\sqrt{3^2 + 4^2} = \qty{5}{cm}$ par an,
dans la direction située à \ang{37} à l'est du nord ($\tan\beta = 3/4$).
(2) $\vec v_{C/A}
= -\vec v_{A/C}$~: \qty{5}{cm} par an dans la direction opposée, à \ang{37}
à l'ouest du sud.
\end{solution}

\begin{solution}{exo:g11:plate-motions:12}
(1) $54/0.375 = \qty{144}{km}$ par million d'années, \qty{14.4}{cm} par an.
(2) De $29/0.375 \approx 77$ à $79/0.375 \approx \qty{211}{km}$ par million
d'années~: entre environ 8 et \qty{21}{cm} par an. (3) Près du point chaud,
la distance n'est pas plus grande que la taille d'un volcan ou de la région
dans laquelle le point chaud alimente des volcans, et l'âge n'est pas plus
grand que le temps que met un volcan à se construire~: tous deux sont très
incertains en valeur relative. Loin du point chaud, une erreur de
\qty{25}{km} sur \qty{2400}{km}, ou d'un demi-million d'années sur
28 millions, ne fait qu'un ou deux pour cent.
\end{solution}

\begin{solution}{exo:g11:plate-motions:13}
(1) Dix fois la dispersion font \qty{50}{mm}~: $50/72 \approx 0.7$ an,
environ huit mois. (2) $50/5 = 10$ ans. (3) En millions d'années, les
plaques se déplacent de dizaines à centaines de kilomètres, des distances
que les bandes magnétiques enregistrent et qu'il était facile de mesurer en
mer~; un positionnement précis à quelques millimètres, qui rend un mouvement
visible en quelques mois, est venu bien plus tard.
\end{solution}

\begin{solution}{exo:g11:plate-motions:14}
(1) En millimètres par an,
\[
  \vec v_{\mathrm{SA/AF}} = (-4.9 - 27.1,\ 11.0 - 17.9) = (-32.0, -6.9).
\] La vitesse vaut $\sqrt{32.0^2 + 6.9^2} \approx
\qty{32.7}{mm}$ par an, et $\tan\alpha = 6.9/32.0 \approx 0.22$ donne
\ang{12} au sud de l'ouest. (2) La \omterm{def:g11:drift-to-spreading:spreading-rate}{vitesse d'expansion totale} est d'environ
\qty{3.3}{cm} par an, la demi-vitesse d'environ \qty{1.6}{cm} par an. (3) La
demi-vitesse des derniers millions d'années, \qty{1.6}{cm} par an, est un
peu plus faible que les \qty{1.9}{cm} par an en moyenne sur 66 millions
d'années~: l'Atlantique Sud s'est peut-être ouvert récemment un peu plus
lentement qu'en moyenne, ou la différence vient peut-être des incertitudes
sur les âges des sédiments. Les deux mesures s'accordent à dire que la
dorsale est lente.
\end{solution}

\begin{solution}{exo:g11:plate-motions:15}
(1) Une droite qui monte de \qty{20}{mm} par an, une chute brusque de
\qty{50}{mm} au bout de cinq ans, puis une droite qui remonte avec la même
pente. (2) La station finit à $10 \times 20 - 50 = \qty{150}{mm}$ à l'est de
son point de départ~: $150/10 = \qty{15}{mm}$ par an, un quart de trop peu.
(3) Chaque moitié donne \qty{20}{mm} par an. (4) Le saut est un mouvement
distinct et soudain causé par le séisme~: le géodésien mesure son ampleur à
partir des positions juste avant et juste après lui, le retire de la série
(ou ajuste les deux parties avec la même pente et une marche entre elles),
et garde \qty{20}{mm} par an comme vitesse de la station.
\end{solution}

\begin{solution}{pb:g11:plate-motions:1}
\textbf{1.} Le Mauna Kea (\num{0.375} million d'années) est le plus jeune,
Suiko (centre) (\num{64.7} millions d'années) le plus ancien. Les âges
augmentent avec la distance à Kilauea.

\textbf{2.} $\qty{1}{km} = 10^6~\unit{mm}$ et un million d'années vaut $10^6$
ans~: \qty{1}{km} par million d'années vaut $10^6~\unit{mm}/10^6$ ans, soit
\qty{1}{mm} par an.

\textbf{3.} $519/5.1 \approx \qty{102}{mm}$ par an.

\textbf{4.} Necker~: $1058/10.3 \approx 102.7$~; Laysan~: $1818/19.9 \approx
91.4$~; Midway~: $2432/27.7 \approx 87.8$~; Colahan~: $3128/38.6 \approx 81.0$~;
Abbott~: $3280/38.7 \approx 84.8$ (millimètres par an).

\textbf{5.} Les points (54~; 0.375), (519~; 5.1), (780~; 7.2), (1058~; 10.3),
(1818~; 19.9), (2432~; 27.7), (3128~; 38.6), (3280~; 38.7) sont proches d'une
droite passant par l'origine, avec une certaine dispersion (par rapport à
une droite de \qty{91}{km} par million d'années, Kauai, Nihoa et Necker se
trouvent un peu en dessous, Laysan dessus, Colahan et Abbott un peu
au-dessus).

\textbf{6.} Chaque volcan s'est construit au-dessus du point chaud et s'est
déplacé avec la plaque depuis. Pendant la durée $t$, la plaque s'est
déplacée de $v\,t$ dans une direction constante~: le volcan se trouve donc à
$d = v\,t$ du point situé maintenant au-dessus du point chaud, et $t = d/v$.

\textbf{7.} $(102.7 + 91.4 + 87.8 + 81.0 + 84.8)/5 \approx \qty{89.5}{mm}$
par an.

\textbf{8.} 89.5 contre \qty{91}{km} par million d'années~: ils diffèrent de
moins de \qty{2}{\percent}.

\textbf{9.} $519/91 \approx 5.7$ millions d'années, pour un âge de
\num{5.1} millions d'années~: c'est proche, Kauai étant un peu plus jeune
que ne le prévoit la droite moyenne.

\textbf{10.} Les volcans sont de plus en plus anciens vers
l'ouest-nord-ouest~: la plaque les emporte dans cette direction. La plaque
Pacifique se déplace vers l'ouest-nord-ouest au-dessus du point chaud.

\textbf{11.} $(4860 - 3493)/(64.7 - 42.4) = 1367/22.3 \approx \qty{61}{km}$
par million d'années, environ \qty{6.1}{cm} par an.

\textbf{12.} $65/91 \approx 0.71$~: le long de la chaîne de l'Empereur, la
plaque se déplaçait environ \qty{30}{\percent} plus lentement que le long de
la chaîne hawaïenne.

\textbf{13.} $3668/50 \approx \qty{73}{km}$ par million d'années, environ
\qty{7.3}{cm} par an, moins que \qty{91}{km} par million d'années. La droite
passant par l'origine utilise les âges de 1987, qui placent la fin de la
chaîne hawaïenne entre 39 et 43 millions d'années~; l'âge plus récent
d'environ 50 millions d'années pour le début du coude étale la même
distance sur une durée plus longue. Les deux sont justes si la plaque s'est
déplacée plus lentement au début de la chaîne hawaïenne et plus vite
ensuite~: la vitesse n'était pas constante.

\textbf{14.} Un changement de la direction du mouvement de la plaque
Pacifique, il y a environ 50 millions d'années.

\textbf{15.} La direction de la chaîne est la direction dans laquelle la
plaque s'est déplacée~; l'espacement des âges le long de la chaîne donne sa
vitesse. Une plaque peut accélérer ou ralentir sans tourner, ou tourner sans
changer de vitesse~: les deux sont donc indépendantes. Des changements de
vitesse sans changement de direction le long de la chaîne de l'Empereur
montrent que l'enregistrement de sa vitesse et le coude sont des événements
distincts~: le coude n'a pas été produit par un ralentissement lié au point
chaud.

\textbf{16.} $\sqrt{63.2^2 + 34.5^2} \approx \qty{72.0}{mm}$ par an, et
$\tan\alpha = 34.5/63.2 \approx 0.55$ donne \ang{29} au nord de l'ouest.

\textbf{17.} L'ouest-nord-ouest est à \ang{22.5} au nord de l'ouest, le
nord-ouest à \ang{45}~: la direction de la station, à \ang{29} au nord de
l'ouest, se situe entre les deux, près de la direction de la chaîne. Elle
est cohérente.

\textbf{18.} La vitesse de la station est mesurée par rapport à un
référentiel dans lequel la \omterm{def:g11:lithosphere-asthenosphere:lithosphere}{lithosphère} dans son ensemble ne tourne pas,
alors que la chaîne donne le mouvement par rapport au point chaud~: si les
deux référentiels se déplacent l'un par rapport à l'autre, les vitesses
diffèrent. Et la station mesure les 30 dernières années, la chaîne les
40 derniers millions d'années, pendant lesquels la vitesse a pu changer. La
dispersion des âges laisse aussi une incertitude de quelques millimètres par
an sur la vitesse de la chaîne.

\textbf{19.} $100/91 \approx 1.1$ million d'années.

\textbf{20.} \textbf{La plaque Pacifique s'est déplacée au-dessus du point
chaud d'Hawaï à environ \qty{9}{cm} par an} pendant les 40 derniers millions
d'années, vers l'ouest-nord-ouest~; avant le coude, elle se déplaçait plus
lentement, et dans une autre direction.
\end{solution}
